\documentclass[11pt]{article}

\usepackage[margin=1in]{geometry}
\usepackage[T1]{fontenc}
\usepackage[utf8]{inputenc}
\usepackage{lmodern}
\usepackage{microtype}
\usepackage{amsmath,amssymb,amsthm,mathtools}
\usepackage{bm}
\usepackage{booktabs}
\usepackage{array}
\usepackage{enumitem}
\usepackage[numbers,sort&compress]{natbib}
\usepackage[colorlinks=true,linkcolor=blue!60!black,citecolor=blue!60!black,urlcolor=blue!60!black]{hyperref}

\newtheorem{definition}{Definition}[section]
\newtheorem{proposition}[definition]{Proposition}
\newtheorem{lemma}[definition]{Lemma}
\newtheorem{corollary}[definition]{Corollary}
\theoremstyle{remark}
\newtheorem{remark}[definition]{Remark}
\newtheorem{example}[definition]{Example}

\newcommand{\R}{\mathbb{R}}
\newcommand{\C}{\mathbb{C}}
\newcommand{\HH}{\mathbb{H}}
\newcommand{\SU}{\mathrm{SU}}
\newcommand{\U}{\mathrm{U}}
\newcommand{\Id}{I}
\newcommand{\ii}{\mathbf{i}}
\newcommand{\jj}{\mathbf{j}}
\newcommand{\kk}{\mathbf{k}}
\newcommand{\norm}[1]{\left\lVert #1 \right\rVert}
\newcommand{\set}[1]{\left\{ #1 \right\}}
\newcommand{\uoneq}{\mathrm{u1q}}
\newcommand{\normalize}{\operatorname{Norm}}
\newcommand{\sgncanon}{\operatorname{SignCan}}
\newcommand{\canon}{\operatorname{Can}}
\newcommand{\Verify}{\operatorname{Verify}}
\newcommand{\MatrixEq}{\operatorname{MatrixEq}}
\newcommand{\TraceDiag}{\operatorname{TraceDiag}}
\newcommand{\ProcFid}{\operatorname{ProcFid}}
\newcommand{\PhaseRes}{\operatorname{PhaseRes}}

\title{\textbf{Verification, Canonicalization, and Equivalence}\\
in Quaternionic Quantum Compilation}

\author{John G. Van Geem\\
RQM Technologies\\
orcid: 0009-0002-4003-8452}

\date{Date 2026-04-09}

\begin{document}
\maketitle

\begin{abstract}
This paper addresses the trust problem for quaternionic quantum compilation. The
preceding papers in the series established a fixed quaternion--$\SU(2)$
convention, a canonical single-qubit intermediate representation
$\uoneq(w,x,y,z)$, and a local fusion pass for single-qubit circuit segments.
The present paper defines how optimized outputs are verified. The core proposal
is an authoritative canonical equivalence criterion: two single-qubit segments
are treated as equivalent when their canonical quaternionic representatives agree
after normalization and sign-canonicalization on the quotient $q\sim -q$
induced by global-phase removal. Around that criterion, the paper defines
preserved invariants, validation procedures, and interoperability diagnostics
with matrix-based toolchains. In particular, it introduces a phase-aware
matrix residual, a trace-overlap diagnostic, and a process-fidelity-style
check under the fixed quaternion--$\SU(2)$ map. A small shipped verification
harness evaluates both positive and negative cases. On 7 measured comparison
cases, including exact named-gate recovery, exact identity elimination,
generic-\uoneq verification, wrong-axis substitutions, and perturbed generic
outputs, all 4 positive cases pass the authoritative canonical test and all
3 negative cases are rejected. The paper argues that quaternionic compilation
can therefore be audited by deterministic canonical output forms while remaining
interoperable with standard matrix-based verification workflows.
\end{abstract}

\noindent\textbf{Keywords:} quantum compilation verification; canonical equivalence; single-qubit segments; $\SU(2)$; unit quaternions; process fidelity diagnostics

\section{Introduction}

Optimization invites a trust question. Once a compiler replaces one gate list by
another, how does one know that the transformation is correct? For quaternionic
compilation the question is especially predictable, because the optimization
language is not the familiar matrix syntax of most toolchains. Papers 2 and 3
in this series showed that single-qubit segments admit a deterministic
quaternionic internal form and that adjacent one-qubit runs can be fused by
quaternion multiplication. The present paper asks the verification question that
follows immediately from those results.

The answer proposed here has two layers. The first layer is \emph{authoritative}
inside the quaternionic compiler: compare canonical representatives. If two
optimized segments canonicalize to the same quaternionic payload, then they are
equivalent modulo global phase under the fixed convention of the series. The
second layer is \emph{interoperable}: expose matrix-based diagnostics that can be
checked by standard toolchains, such as phase-aware residuals, trace overlap,
and process fidelity under the same $\SU(2)$ representative.

The point is not to replace matrix-based verification. The point is to make
quaternionic compilation auditable in its own native representation while
remaining transparent to partner systems, SDKs, and review workflows that still
speak in matrices. This paper therefore treats verification as part of the
compiler method itself.

Its scope is deliberately narrow. It concerns single-qubit segments only. It
does not define verification of entangling-gate internals, does not claim a
universal proof framework for all quantum optimizations, and does not claim that
one diagnostic alone is sufficient for every integration setting.

\section{Background dependency on Quaternionic Quantum Computing}

First we fix the quaternion--$\SU(2)$ convention:
\begin{equation}
\Phi(a+b\ii+c\jj+d\kk)
=
\begin{pmatrix}
a-di & -c-bi \\
c-bi & a+di
\end{pmatrix}
=
a\Id-i(b\sigma_x+c\sigma_y+d\sigma_z).
\label{eq:phi}
\end{equation}
Under this map, the unit quaternions model $\SU(2)$ exactly, and every
single-qubit gate action modulo global phase is represented by a sign class
$\set{q,-q}$ of unit quaternions.

Second we define the canonical IR object
\[
\uoneq(w,x,y,z)
\]
with the canonicalization rule
\[
\canon=\sgncanon\circ\normalize,
\]
where $\normalize$ restores unit norm and $\sgncanon$ chooses the preferred
representative by the rule $w>0$, with an explicit deterministic tie-break on
the equator $w=0$.

Third we define the local fusion operator for single-qubit segments
\[
\Fuse(\mathrm{seg})=\canon(q_n\cdots q_1)
\]
and used canonical quaternion equality as the authoritative equivalence criterion
for fused outputs. The present paper turns that criterion into a formal
verification layer.

\section{Problem setup: the trust layer}

\begin{definition}[Single-qubit segment]
A \emph{single-qubit segment} on wire $a$ is a maximal consecutive run of gates
that act only on wire $a$, bounded by the beginning or end of the circuit, or
by a boundary operation outside the present single-qubit model for that wire,
including entangling gates, measurement, reset, barrier, or classical-control
boundaries.
\end{definition}

\begin{definition}[Verification problem]
Given an original single-qubit segment $\mathrm{seg}$ and a candidate optimized
segment $\mathrm{seg}'$, the verification problem is to determine whether
$\mathrm{seg}$ and $\mathrm{seg}'$ represent the same single-qubit action modulo
global phase under the fixed quaternion--$\SU(2)$ convention.
\end{definition}

\begin{remark}
This paper treats verification as a relation between \emph{segment actions}, not
between textual syntaxes. Two segments may use different named-gate spellings,
different reconstruction policies, or a generic $\uoneq$ primitive, yet still be
equivalent at the represented-action level.
\end{remark}

\section{Canonical equivalence and preserved invariants}

\subsection{Authoritative canonical criterion}

\begin{definition}[Canonical segment representative]
Let $\mathrm{seg}=(g_1,\dots,g_n)$ be a single-qubit segment in application
order, and let $q_1,\dots,q_n$ be the translated unit-quaternion representatives
of those gates. The canonical segment representative is
\[
Q(\mathrm{seg})=\canon(q_n\cdots q_1).
\]
\end{definition}

\begin{definition}[Authoritative equivalence]
Two single-qubit segments $\mathrm{seg}$ and $\mathrm{seg}'$ are
\emph{canonically equivalent}, written
\[
\Verify(\mathrm{seg},\mathrm{seg}'),
\]
if
\[
Q(\mathrm{seg})=Q(\mathrm{seg}').
\]
\end{definition}

\begin{proposition}[Canonical equivalence is exact segment equivalence modulo global phase]
For single-qubit segments, $\Verify(\mathrm{seg},\mathrm{seg}')$ holds if and
only if $\mathrm{seg}$ and $\mathrm{seg}'$ represent the same $\SU(2)$ action
up to the global-phase quotient inherited from $\U(2)$.
\end{proposition}

\begin{proof}
By paper 1, single-qubit actions modulo global phase correspond exactly to sign
classes $\set{q,-q}$ of unit quaternions under the fixed map \eqref{eq:phi}. By
paper 2, $\canon$ chooses exactly one representative from each sign class. Hence
two segments produce the same canonical representative if and only if they
produce the same sign class, which is exactly the same single-qubit action modulo
global phase.
\end{proof}

\subsection{Preserved invariants}

\begin{definition}[Quaternionic compilation invariants]
For a valid optimized replacement of a single-qubit segment, the following
invariants are preserved:
\begin{enumerate}[label=(\alph*), leftmargin=1.5em]
\item unit norm after normalization;
\item canonical sign choice under the fixed hemisphere rule;
\item represented $\SU(2)$ action modulo global phase; and
\item reproducible canonical output form.
\end{enumerate}
\end{definition}

\begin{remark}
The final item is operationally important. Even when multiple frontend or
presentation-level reconstructions are available, the canonical quaternionic
payload is unique. This is what makes optimized outputs reproducible across runs
and comparable across implementations.
\end{remark}

\begin{proposition}[Canonical output form is deterministic]
If two compilation runs produce valid optimized representations of the same
single-qubit segment action under the fixed convention, then their canonical
quaternionic payloads are identical.
\end{proposition}

\begin{proof}
Both outputs represent the same sign class $\set{q,-q}$ in the unit quaternions.
The canonicalization map of paper 2 chooses exactly one representative from that
class. Therefore the payload is unique.
\end{proof}

\section{Validation procedures}

\subsection{Primary validation procedure}

\begin{definition}[Canonical validation procedure]
To validate a candidate optimized segment $\mathrm{seg}'$ against an original
segment $\mathrm{seg}$:
\begin{enumerate}[label=(\arabic*), leftmargin=1.5em]
\item translate each segment into the fixed quaternionic representation;
\item multiply in application order to obtain raw segment products;
\item normalize and sign-canonicalize both products; and
\item compare the resulting canonical representatives for equality.
\end{enumerate}
\end{definition}

\begin{remark}
This procedure is the authoritative one for the present series. It is exact in
the mathematical model and becomes an implementation check after explicit
tolerance rules are layered on top.
\end{remark}

\subsection{Tolerance-aware numerical variant}

\begin{definition}[Tolerance-aware canonical check]
In floating-point settings, treat two canonical representatives as equal when
their coordinatewise difference is below a declared tolerance after
normalization and sign-canonicalization.
\end{definition}

\begin{remark}
The tolerance belongs to the implementation policy, not to the exact
equivalence relation. The exact mathematical criterion remains equality of
canonical representatives.
\end{remark}

\section{Interoperability with matrix-based toolchains}

Trust often requires compatibility with existing verification culture. For this
reason the paper introduces diagnostics that can be computed from the same
segment pair after applying the fixed map \eqref{eq:phi}.

\begin{definition}[Phase-aware matrix equivalence]
Let $U,V\in\SU(2)$ be the matrix representatives of two single-qubit segments.
Define
\[
\MatrixEq(U,V)
\]
to hold when $U^\dagger V$ is equal to a scalar phase times the identity.
\end{definition}

\begin{definition}[Trace overlap and process fidelity]
For $U,V\in\SU(2)$ define the trace-overlap diagnostic
\[
\TraceDiag(U,V)=\left|\mathrm{Tr}(U^\dagger V)\right|
\]
and the process-fidelity-style quantity
\[
\ProcFid(U,V)=\frac{\left|\mathrm{Tr}(U^\dagger V)\right|^2}{4}.
\]
\end{definition}

\begin{definition}[Phase-aware residual]
Let $\lambda=\mathrm{Tr}(U^\dagger V)/|\mathrm{Tr}(U^\dagger V)|$ when the trace
is nonzero. Define the phase-aware residual
\[
\PhaseRes(U,V)
\]
as the maximum-entry norm of $U^\dagger V-\lambda \Id$.
\end{definition}

\begin{proposition}[Canonical equality implies perfect matrix diagnostics]
If $\Verify(\mathrm{seg},\mathrm{seg}')$ holds exactly, then for the
corresponding matrix representatives $U$ and $V$ one has
\[
\MatrixEq(U,V),
\qquad
\ProcFid(U,V)=1,
\qquad
\TraceDiag(U,V)=2,
\qquad
\PhaseRes(U,V)=0.
\]
\end{proposition}

\begin{proof}
Canonical equality implies the same unit quaternion representative. Under
\eqref{eq:phi}, the two segments therefore produce the same $\SU(2)$ matrix, so
$U^\dagger V=\Id$. The stated diagnostics follow immediately.
\end{proof}

\begin{remark}
The converse must be handled with care numerically. Near-equality in trace or
fidelity can diagnose closeness, but the authoritative exact yes/no criterion in
this paper remains canonical quaternion equality after canonicalization.
\end{remark}

\section{Worked examples}

\begin{example}[Exact named-gate recovery]
Compare
\[
R_x(\pi/2);R_x(\pi/2)
\qquad\text{and}\qquad
X.
\]
Both sides canonicalize to the quaternion $(0,1,0,0)$. Hence
$\Verify(\mathrm{seg},\mathrm{seg}')$ holds exactly.
\end{example}

\begin{example}[Identity elimination]
Compare
\[
T;T^\dagger
\qquad\text{and}\qquad
\emptyset.
\]
Both sides canonicalize to $(1,0,0,0)$, so the empty output is verified as a
correct identity elimination.
\end{example}

\begin{example}[Generic primitive verification]
Let the original segment be
\[
H;T;R_z(\pi/3),
\]
and let the optimized candidate be the generic primitive
\[
\uoneq(0.560985526797,-0.430459334577,-0.560985526797,-0.430459334577).
\]
The two canonical quaternionic payloads agree, so the generic primitive is
verified exactly even though it is not a named gate.
\end{example}

\begin{example}[Negative wrong-axis case]
Compare
\[
R_z(\pi/3)
\qquad\text{and}\qquad
R_x(\pi/3).
\]
The canonical representatives differ, so the authoritative canonical check fails.
The matrix diagnostics also reflect non-equivalence:
\[
\ProcFid=0.5625,\qquad \PhaseRes=0.5.
\]
\end{example}

\section{Reference verification corpus and measured diagnostics}

The paper package includes the executable verifier
\texttt{artifacts/code/verify\_canonical\_equivalence.py}, which generates the
measured results stored in \texttt{artifacts/code/verification\_results.json}.
The shipped reference corpus contains 7 comparison cases:
\begin{itemize}[leftmargin=1.5em]
\item 4 positive equivalence cases; and
\item 3 negative non-equivalence cases.
\end{itemize}

The measured aggregate outcome is:
\begin{itemize}[leftmargin=1.5em]
\item all 4 positive cases pass the authoritative canonical check;
\item all 3 negative cases are rejected by the authoritative canonical check; and
\item all 7 cases are consistent with the expected verdict.
\end{itemize}

Table~\ref{tab:verification-results} summarizes the reference corpus.

\begin{table}[h]
\centering
\begin{tabular}{@{}lcccc@{}}
\toprule
Case type & Cases & Expected equivalent & Passed/rejected & Consistent \\
\midrule
Positive exact cases & 4 & yes & 4 passed & yes \\
Negative distinct cases & 3 & no & 3 rejected & yes \\
\midrule
Total & 7 & mixed & 7 correct & yes \\
\bottomrule
\end{tabular}
\caption{Measured verification results on the shipped reference corpus.}
\label{tab:verification-results}
\end{table}

The individual diagnostic patterns are also informative:
\begin{enumerate}[label=(\arabic*), leftmargin=1.5em]
\item positive exact cases produce canonical equality, process fidelity $1$, and
near-zero phase-aware residuals;
\item structurally wrong substitutions such as $X$ versus $Y$ produce zero
trace overlap and zero process fidelity; and
\item small perturbations of a generic $\uoneq$ primitive can exhibit fidelity
close to $1$ while still failing the authoritative canonical equality test.
\end{enumerate}

\begin{remark}
The third point is operationally important. Fidelity-like diagnostics are useful
for interoperability and debugging, but they are not by themselves a substitute
for exact canonical equivalence when the compiler intends an exact optimization.
\end{remark}

\section{Standard mathematics vs.\ RQM Technologies framing}

\subsection{Standard ingredients}

The following ingredients are standard:
\begin{enumerate}[label=(\arabic*), leftmargin=1.5em]
\item single-qubit gate actions modulo global phase are modeled by $\SU(2)$
representatives;
\item unit quaternions model $\SU(2)$ under the fixed convention of paper 1;
\item the ambiguity in that model is exactly $q\sim -q$; and
\item matrix overlap, trace diagnostics, and fidelity-style checks are standard
verification tools.
\end{enumerate}

\subsection{Project-specific contribution}

The project-specific contribution of this paper is the trust layer:
\begin{enumerate}[label=(\alph*), leftmargin=1.5em]
\item canonical quaternion equality as the authoritative exact verification rule;
\item fixed sign conventions and reproducible canonical output forms;
\item explicit preserved invariants for quaternionic compilation; and
\item a paired matrix-interoperability layer for integration with existing
toolchains.
\end{enumerate}

\section{Scope, limitations, and non-claims}

This paper does \emph{not}:
\begin{enumerate}[label=(\arabic*), leftmargin=1.5em]
\item define verification of entangling-gate internals;
\item provide a universal proof system for all quantum compiler passes;
\item claim that trace or fidelity diagnostics alone are sufficient for every
integration context;
\item prove anything about hardware execution correctness; or
\item replace existing matrix-based verification practice.
\end{enumerate}

It defines a trust layer for single-qubit quaternionic compilation: canonical
equivalence, preserved invariants, validation procedures, and interoperable
diagnostics.

\section{Conclusion}

Quaternionic compilation becomes safer to adopt once its outputs can be verified
deterministically. The central proposal of this paper is therefore simple:
\begin{enumerate}[label=(\arabic*), leftmargin=1.5em]
\item canonical quaternion equality is the authoritative exact equivalence
criterion for single-qubit segments;
\item preserved invariants make optimized outputs reproducible; and
\item matrix-based diagnostics provide interoperability with standard
toolchains.
\end{enumerate}

In that sense, this paper supplies the trust layer for the earlier
representation and optimization papers. It does not change the compiler method;
it makes that method auditable.

\appendix

\section{Recommended validation pipeline}

For an exact local optimization on a single-qubit segment, the recommended
verification procedure is:
\begin{enumerate}[label=(\arabic*), leftmargin=1.5em]
\item compute canonical quaternionic representatives for original and optimized
segments;
\item compare them for authoritative exact equality;
\item if needed for interoperability, compute the matrix representatives under
\eqref{eq:phi};
\item record phase-aware residual, trace overlap, and process fidelity for audit
logs.
\end{enumerate}

\section{Companion code and figures}

The executable verifier and its measured output are included in
\texttt{artifacts/code/verify\_canonical\_equivalence.py} and
\texttt{artifacts/code/verification\_results.json}. Publication-oriented SVG
figures are included in \texttt{artifacts/figures/}.

\bibliographystyle{plainnat}
\bibliography{references}

\end{document}
